\section{Introduction}
\label{sec:introduction}

A miss-ratio curve (MRC) gives the hit ratio $H(C)$ of a cache of capacity $C$
on a given request stream. Operators use it to size memory: the capacity at
which the curve flattens is the point past which additional memory buys little
hit ratio. Published fleet analyses show that capacity decisions of this kind
are made continuously across hundreds of clusters \cite{yang2020}.

\subsection{The Gap: Aggregate Regression}

The dominant practical approach regresses an aggregate hit count, or an
aggregate miss ratio, on capacity. The fit is performed over a set of measured
capacities, and the curve is interpolated or extrapolated from those points.
This formulation has two costs. First, capacity enters as a covariate, so every
new capacity requires new measurements or a new fit. Second, aggregating over
requests removes the per-request quantity that actually governs the curve: how
far back in the stream a request must reach to find its previous access. Two
workloads with identical aggregate hit counts at one capacity can have
different curves at every other capacity, and aggregate regression cannot
distinguish them.

\subsection{Motivation: Over- and Under-Allocation}

Both directions of sizing error are expensive. Over-allocation reserves memory
that produces no hit-ratio gain, which in a multi-tenant fleet is memory taken
from another workload. Under-allocation produces misses that fall through to
the storage tier, where the latency penalty is orders of magnitude larger than
a cache lookup. A curve that is accurate in shape, and not only at the
capacities that were measured, supports sizing at capacities that were never
observed. That is the operational motivation for modelling the per-request
reuse process instead of its aggregate.

\subsection{Contributions}

\begin{enumerate}[nosep]
    \item A formulation of reuse distance as a right-censored discrete
    time-to-event variable, with the LRU stack property supplying the exact map
    from its survival function to the MRC (Section~\ref{sec:formalization}).
    \item Two gradient-boosted discrete hazard parametrizations, one treating
    reuse distance as a feature and one fitting a separate model per hazard
    bin, together with an explicit censoring correction and a documented
    Optuna search space (Section~\ref{sec:methods}).
    \item An evaluation protocol against simulation ground truth on production
    traces, reporting MRC error, hazard calibration, and the integrated Brier
    score rather than a single aggregate fit statistic
    (Section~\ref{sec:experiments}).
\end{enumerate}

\subsection*{Scope and Page Budget}
\phantomsection\label{sec:pagebudget}

% PAGE BUDGET: body 8-10 pages (Sections 1-6). References and appendix are
% outside the budget. Derivations, full hyperparameter tables, per-cluster
% numbers and reproduction commands live in the appendix by design, so that
% the body states the argument and the appendix carries the evidence.
The body of this report is budgeted at eight to ten pages, following the
allocation recorded in \texttt{main.tex} and anchored at
Section~\ref{sec:pagebudget}. Material that does not fit moves to the appendix
rather than being compressed. The report covers LRU-referenced MRCs only;
non-LRU replacement is discussed as a limitation in
Section~\ref{sec:conclusion}.
